% ===== 06_joint_control =====
\section{接触率与隔离率的联合阈值控制}
\label{sec:joint}
固定接触率时，维持 $I=\eta$ 唯一确定隔离率。由推论\ref{cor:s1:qc-structure}，隔离能力不足时，接触削减可能补足实施能力；仅隔离已经可行时，联合控制仍可用于研究不同干预分配的成本。现在允许接触率和隔离率同时变化，感染人数不变只约束二者的乘积，因而还需要确定如何分配两种干预。本节先给出全部联合阈值控制及其可行条件，再比较持续时间、累计感染和隔离人数，最后确定给定成本下的类内最优分配。

以下固定由常规轨迹确定的 $t_1,S^*$ 及阈值 $\eta$，其中 $S^*>S_c$。所比较的控制均从 $(S^*,\eta)$ 出发，在整个控制期保持 $I=\eta$，首次达到 $S=S_c$ 后恢复 $(c_0,q_0)$。控制有界可测，满足 $0<c\le c_0$、$q_0\le q\le1$，状态取绝对连续解；几乎处处相等的控制视为同一控制。这里固定的是进入与退出时的 $(S,I)$，其余仓室在退出时可以不同。控制持续时间 $\Delta t=t_2-t_1$ 随所选策略改变。

\subsection{联合控制的状态表达与可行条件}
由于控制期内 $S$ 严格下降，可以用 $S$ 代替时间描述控制。沿用前文 $q_c(S)$ 的记法，分别以 $c_c(S),q_c(S)$ 表示接触率和隔离率的状态表达；它们对应的时间函数为 $c(t)=c_c(S(t))$、$q(t)=q_c(S(t))$。仅隔离控制就是 $c_c(S)\equiv c_0$ 的特例。
\begin{theorem}
\label{thm:joint:representation}
规定的联合阈值控制与满足
\begin{equation}
\frac{c_0S_c}{S}\le c_c(S)\le c_0,\qquad S\in[S_c,S^*],
\label{eq:joint:contact-range}
\end{equation}
的有界可测函数相对应。隔离率由
\begin{equation}
q_c(S)=1-\frac{\gamma N}{\beta c_c(S)S}
      =1-\frac{c_0(1-q_0)S_c}{c_c(S)S}
\label{eq:joint:q}
\end{equation}
确定，状态轨迹与控制持续时间满足
\begin{align}
t-t_1&=\frac{N}{\eta}\int_{S(t)}^{S^*}
\frac{du}{c_c(u)u-c_0\bar S},\label{eq:joint:time}\\
\Delta t[c_c]&=\frac{N}{\eta}\int_{S_c}^{S^*}
\frac{dS}{c_c(S)S-c_0\bar S}.\label{eq:joint:duration}
\end{align}
式\eqref{eq:joint:time}唯一确定严格递减的 $S(t)$。选取满足式\eqref{eq:joint:contact-range}的逐点代表后，退出时 $c(t)\to c_0$、$q(t)\to q_0$；若 $c_c(S)$ 连续，则两控制在控制期内部连续。
\end{theorem}
\begin{proof}
在 $I=\eta>0$ 上，感染者方程给出
\begin{equation}
c(t)(1-q(t))S(t)=\frac{\gamma N}{\beta}
=c_0(1-q_0)S_c.
\label{eq:joint:balance}
\end{equation}
因而正的恒定感染阶段不允许 $q=1$。由 $q\ge q_0$ 和 $c\le c_0$，可得式\eqref{eq:joint:contact-range}和\eqref{eq:joint:q}；反向代入也成立。利用式\eqref{eq:joint:balance}，易感者方程化为
\begin{equation}
\dot S=-\frac{\eta}{N}\bigl[c_c(S)S-c_0\bar S\bigr].
\label{eq:joint:Sdot}
\end{equation}
其下降速度满足
\[
0<c_0(S_c-\bar S)\le c_c(S)S-c_0\bar S
\le c_0(S^*-\bar S).
\]
因此 $S$ 在有限时间从 $S^*$ 到达 $S_c$，分离变量即得式\eqref{eq:joint:time}和\eqref{eq:joint:duration}。

反之，给定满足式\eqref{eq:joint:contact-range}的可测函数，可在零测集上修改使界限处处成立。式\eqref{eq:joint:time}右侧随积分下限严格递减，且该时间变换及其逆函数均为 Lipschitz 函数，故定义唯一绝对连续的 $S(t)$。用式\eqref{eq:joint:q}构造时间控制，令 $I=\eta$，则前两个状态方程几乎处处成立；其余仓室由原模型确定。该时间变换保持零测集，所以这一对应在几乎处处相等意义下成立。最后，由 $c_0S_c/S\le c_c(S)\le c_0$，当 $S\downarrow S_c$ 时两界均趋于 $c_0$；式\eqref{eq:joint:q}继而给出 $q\to q_0$。
\end{proof}

这一表示只需选择接触率 $c_c(S)$，隔离率和控制时间随后确定。由此，两个时间函数之间的约束化为接触率状态函数的取值范围。

进一步设最低允许接触率为 $c_{\min}$，最高可达到的隔离比例为 $q_{\rm cap}$，其中
\[
0\le c_{\min}\le c_0,\qquad q_0\le q_{\rm cap}<1.
\]
$c_{\min}$ 是实施能力限制，与前文用于判断是否触发控制的 $c_{0,\min}(\eta)$ 不同；$q_{\rm cap}$ 也不同于仅隔离策略实际需要的启动最大值 $q_{\max}$。为简写取值范围，记
\begin{equation}
\underline c(S)=\max\!\left\{c_{\min},\frac{c_0S_c}{S}\right\},\qquad
\overline c(S)=\min\!\left\{c_0,
\frac{c_0(1-q_0)S_c}{(1-q_{\rm cap})S}\right\}.
\label{eq:joint:bounds}
\end{equation}
\begin{corollaryn}
\label{cor:joint:feasibility}
上述能力限制等价于 $\underline c(S)\le c_c(S)\le\overline c(S)$。从 $S^*$ 至 $S_c$ 的整段阈值控制可行，当且仅当
\begin{equation}
\frac{\beta c_{\min}(1-q_{\rm cap})S^*}{N}\le\gamma.
\label{eq:joint:feasibility}
\end{equation}
利用式\eqref{eq:s1:qmax}，该条件也可写为
\[
c_{\min}(1-q_{\rm cap})\le c_0(1-q_{\max}).
\]
\end{corollaryn}
\begin{proof}
$c_c\ge c_{\min}$ 与式\eqref{eq:joint:contact-range}给出下界；由式\eqref{eq:joint:q}中的 $q_c\le q_{\rm cap}$ 得到上界。在 $S\in[S_c,S^*]$ 上，已有 $c_0S_c/S\le c_0$、$c_0S_c/S\le c_0(1-q_0)S_c/[(1-q_{\rm cap})S]$ 及 $c_{\min}\le c_0$，故只需检验
\[
c_{\min}\le\frac{c_0(1-q_0)S_c}{(1-q_{\rm cap})S}.
\]
其最严格处为 $S=S^*$，即式\eqref{eq:joint:feasibility}。条件成立时取 $c_c=(\underline c+\overline c)/2$，由前一定理得到可行控制；条件不成立时，启动状态附近不存在允许分配。
\end{proof}
由于 $c_{\min}\le c_0$，联合可行条件不强于仅隔离条件；当 $c_{\min}=c_0$ 时，二者相同。具体而言，若 $q_{\rm cap}\ge q_{\max}$，仅隔离与联合控制均可行；允许分配存在自由度时，仍可比较不同分配的持续时间、累计感染和成本。若 $q_{\rm cap}<q_{\max}$ 而式\eqref{eq:joint:feasibility}成立，仅隔离不可行，但接触削减可以补足隔离能力。沿用满足界限的逐点代表约定，由式\eqref{eq:joint:bounds}可知，在
\[
\frac{\gamma N}{\beta c_0(1-q_{\rm cap})}<S\le S^*
\]
上必须有 $c_c(S)<c_0$，因此接触削减至少在启动后的一段正时间内是必要的；这不要求隔离率也必须高于 $q_0$。当 $S$ 降至上述区间下界及其以下时，仅隔离重新成为允许分配，但不意味着成本最优策略会切回仅隔离。最后，若式\eqref{eq:joint:feasibility}不成立，规定首次触发状态附近不存在允许分配，因而不能等到当前常规轨迹达到 $I=\eta$ 才开始本节控制；提前干预改变启动状态的方案仍需另行分析。Angulo 等\cite{Angulo2021}在 SIR 模型中讨论了能力有限时需要提前启动的情形，但这一结果不直接保证本文 SIQR 模型中的提前干预可行。

由引理\ref{lem:s1:Sstar-sensitivity}，$S^*$ 随 $\eta$ 严格下降。因此，在其余参数及能力界限固定时，只要某一内部阈值能够实施，所有更高且仍会触发控制的阈值也能实施。降低感染上限不仅延长控制时间，还可能使启动时所需干预超过能力限制。

\subsection{一类显式联合控制}
一般的 $c_c(S)$ 由积分确定时间轨迹，未必能以初等函数反演。取分配参数 $\alpha\in[0,1]$，构造
\begin{align}
c_c(S;\alpha)&=c_0\left[\alpha+(1-\alpha)\frac{S_c}{S}\right],\label{eq:joint:explicit-c}\\
q_c(S;\alpha)&=1-\frac{(1-q_0)S_c}{S_c+\alpha(S-S_c)}.
\label{eq:joint:explicit-q}
\end{align}
$\alpha$ 仅表示干预分配；原记号 $\theta=\eta/N$ 仍表示阈值比例，$\lambda$ 仍表示拐点相对位置。式\eqref{eq:joint:explicit-c}满足\eqref{eq:joint:contact-range}，所以它给出全部联合阈值控制中的一个显式子族。
\begin{theorem}
\label{thm:joint:explicit}
固定 $0<\alpha\le1$ 时，式\eqref{eq:joint:explicit-c}--\eqref{eq:joint:explicit-q}对应的控制期状态为
\begin{equation}
S(t)=S_c+
\left(S^*-S_c+\frac{S_c-\bar S}{\alpha}\right)
 e^{-\frac{c_0\eta\alpha}{N}(t-t_1)}
-\frac{S_c-\bar S}{\alpha},\qquad I(t)=\eta,
\label{eq:joint:explicit-S}
\end{equation}
控制持续时间为
\begin{equation}
\Delta t(\alpha)=\frac{N}{c_0\eta\alpha}
\ln\!\left(1+\frac{\alpha(S^*-S_c)}{S_c-\bar S}\right).
\label{eq:joint:explicit-duration}
\end{equation}
$\alpha=0$ 时，相应表达为
\begin{equation}
S(t)=S^*-\frac{c_0\eta}{N}(S_c-\bar S)(t-t_1),\qquad
\Delta t(0)=\frac{N(S^*-S_c)}{c_0\eta(S_c-\bar S)}.
\label{eq:joint:explicit-zero}
\end{equation}
将上述 $S(t)$ 代入式\eqref{eq:joint:explicit-c}--\eqref{eq:joint:explicit-q}，即得到两个控制的时间表达。$\alpha=0$ 对应仅减少接触，$\alpha=1$ 恢复第\ref{sec:s1}节的仅隔离控制。对于 $0<\alpha<1$，控制期内接触率严格增加、隔离率严格下降，并在退出时同时连续恢复为 $(c_0,q_0)$。
\end{theorem}
\begin{proof}
代入式\eqref{eq:joint:Sdot}得
\[
\dot S=-\frac{c_0\eta}{N}\bigl[S_c-\bar S+\alpha(S-S_c)\bigr],
\qquad S(t_1)=S^*.
\]
分别求解 $\alpha>0$ 和 $\alpha=0$ 的线性方程，得到状态表达；令 $S(t_2)=S_c$ 得到持续时间。两组公式在 $\alpha\downarrow0$ 时连续衔接。对于 $0<\alpha<1$，
\[
\frac{\partial c_c(S;\alpha)}{\partial S}
=-\frac{c_0(1-\alpha)S_c}{S^2}<0,\qquad
\frac{\partial q_c(S;\alpha)}{\partial S}
=\frac{(1-q_0)S_c\alpha}{[S_c+\alpha(S-S_c)]^2}>0.
\]
结合 $\dot S<0$，得到时间单调性。代入 $S=S_c$ 得到退出值；代入 $\alpha=0,1$ 分别恢复两种单控制。
\end{proof}
控制启动时，$0<\alpha<1$ 的策略使 $c$ 向下跳跃、$q$ 向上跳跃。由于 $S^*>S_c$ 时常规控制不满足式\eqref{eq:joint:balance}，这一跳跃与退出时的连续恢复不同。
\begin{corollaryn}
\label{cor:joint:allocation-range}
在式\eqref{eq:joint:bounds}的能力限制下，上述显式族可行当且仅当
\begin{equation}
\max\!\left\{0,\frac{c_{\min}S^*-c_0S_c}{c_0(S^*-S_c)}\right\}
\le\alpha\le
\min\!\left\{1,\frac{S_c(q_{\rm cap}-q_0)}{(1-q_{\rm cap})(S^*-S_c)}\right\}.
\label{eq:joint:allocation-range}
\end{equation}
该区间非空当且仅当式\eqref{eq:joint:feasibility}成立。
\end{corollaryn}
\begin{proof}
显式族中 $c$ 不减、$q$ 不增，故只需在 $S=S^*$ 处检查 $c\ge c_{\min}$ 和 $q\le q_{\rm cap}$。代入式\eqref{eq:joint:explicit-c}--\eqref{eq:joint:explicit-q}并与 $\alpha\in[0,1]$ 取交，得到式\eqref{eq:joint:allocation-range}。未截断的下界不超过 $1$，未截断的上界非负；二者的大小关系等价于
$c_{\min}S^*\le c_0(1-q_0)S_c/(1-q_{\rm cap})$。
这正是式\eqref{eq:joint:feasibility}，故该区间与完整控制类同时非空。
\end{proof}
这一子族用于给出可计算的控制和两种单控制之间的连续分配，不代表全部联合控制，也不声称它是唯一可以显式求解的选择。下面的比较与成本最小化均返回全部允许的 $c_c(S)$。

\subsection{控制持续时间、累计感染与隔离人数}
在相同易感者状态上，较大的 $c_c(S)$ 对应较快的 $S$ 下降。因此，当进入与退出状态固定时，持续时间可直接由式\eqref{eq:joint:duration}比较。
\begin{theorem}
\label{thm:joint:comparison}
若两个允许控制的状态函数满足 $c_c^{(1)}(S)\ge c_c^{(2)}(S)$ 几乎处处成立，则
$\Delta t[c_c^{(1)}]\le\Delta t[c_c^{(2)}]$。若前一不等式在正测度集合上严格成立，则后一不等式严格。特别地，
\begin{equation}
\frac{N}{c_0\eta}\ln\frac{S^*-\bar S}{S_c-\bar S}
\le\Delta t[c_c]\le
\frac{N(S^*-S_c)}{c_0\eta(S_c-\bar S)}.
\label{eq:joint:duration-bounds}
\end{equation}
左侧等号当且仅当 $c_c(S)=c_0$ 几乎处处，即仅加强隔离；右侧等号当且仅当 $c_c(S)=c_0S_c/S$ 几乎处处，即仅减少接触。显式族的 $\Delta t(\alpha)$ 在 $[0,1]$ 上严格递减。

沿任一允许控制，其控制期累计新增感染与新增隔离易感者人数分别满足
\begin{align}
\Itcum(t_2)-\Itcum(t_1)
&=\beta(S^*-S_c)+\gamma\eta(1-\beta)\Delta t,
\label{eq:joint:cumulative}\\
S_q(t_2)-S_q(t_1)
&=(1-\beta)\bigl[S^*-S_c-\gamma\eta\Delta t\bigr].
\label{eq:joint:quarantine}
\end{align}
当 $0<\beta<1$ 时，控制期累计新增感染与持续时间具有相同排序，新增隔离易感者人数的排序相反；$\beta=1$ 时，累计新增感染均为 $S^*-S_c$，隔离易感者人数不增加。
\end{theorem}
\begin{proof}
式\eqref{eq:joint:duration}的被积函数关于 $c_c(S)$ 严格递减，故得到一般排序。由式\eqref{eq:joint:contact-range}又有
\[
\frac1{c_0(S-\bar S)}\le
\frac1{c_c(S)S-c_0\bar S}\le
\frac1{c_0(S_c-\bar S)}.
\]
积分给出两界。等号成立要求相应非负积分差为零，即两个状态函数几乎处处相等。显式族中 $c_c(S;\alpha)$ 对 $S>S_c$ 随 $\alpha$ 严格增加，故其持续时间严格减少。

积分式\eqref{eq:joint:Sdot}，得
\[
S^*-S_c=\frac{\eta}{N}\int_{t_1}^{t_2}c(t)S(t)\,dt
-\frac{c_0\eta\bar S}{N}\Delta t.
\]
左式积分项乘以 $\beta$ 即为控制期总新增感染。由 $\beta c_0\bar S/N=\gamma(1-\beta)$，得到式\eqref{eq:joint:cumulative}。控制期离开 $S$ 的人群分为新感染者与未感染的隔离者，故
\[
S^*-S_c=\Itcum(t_2)-\Itcum(t_1)+S_q(t_2)-S_q(t_1),
\]
代入前式得到式\eqref{eq:joint:quarantine}。
\end{proof}
无额外能力限制时，仅隔离控制达到规定控制类内最短的持续时间；$0<\beta<1$ 时，它也达到最少的控制期累计感染，但同时产生最多的新增隔离易感者。这源于未感染者隔离加快了社区易感者的减少，并非说明增加接触本身有利于控制疫情。上述隔离人数也不等于隔离人天数。

加入能力限制后，两种单控制可能不再可行。应改在实际允许范围内比较，而不能继续把不可实施的控制作为可达到的界限。
\begin{corollaryn}
\label{cor:joint:capped-comparison}
若式\eqref{eq:joint:feasibility}成立，则
\begin{equation}
\Delta t[\overline c]\le\Delta t[c_c]\le\Delta t[\underline c],
\label{eq:joint:capped-duration}
\end{equation}
其中两端由式\eqref{eq:joint:duration}分别代入 $\overline c(S)$、$\underline c(S)$ 计算，且均可达到。达到左、右等号分别要求 $c_c=\overline c$、$c_c=\underline c$ 几乎处处成立。$0<\beta<1$ 时，同一两个控制分别达到受限类内最少和最多的控制期累计新增感染。
\end{corollaryn}
\begin{proof}
在式\eqref{eq:joint:duration}中用 $\underline c\le c_c\le\overline c$ 比较即可。两个边界函数连续且满足限制，因而确实对应可行控制。感染量结论由式\eqref{eq:joint:cumulative}得到。
\end{proof}

控制期的排序还能用于比较全过程，但必须保留共同的启动规则和退出状态。
\begin{corollaryn}
\label{cor:joint:whole-epidemic}
设两种允许策略在启动前及退出后均采用 $(c_0,q_0)$，并从同一初始状态首次到达同一阈值 $\eta>1$。则
\begin{align}
t_{\rm end}^{(1)}-t_{\rm end}^{(2)}
&=\Delta t^{(1)}-\Delta t^{(2)},\label{eq:joint:whole-time}\\
\Itcum^{(1)}-\Itcum^{(2)}
&=\gamma\eta(1-\beta)\bigl(\Delta t^{(1)}-\Delta t^{(2)}\bigr).
\label{eq:joint:whole-cumulative}
\end{align}
第二式对计算至无限时刻的累计新增感染同样成立。能力限制下，$\overline c(S)$ 因而同时达到最短控制时间、最早社区动态清零时刻；当 $0<\beta<1$ 时，也达到最少的全过程累计新增感染。
\end{corollaryn}
\begin{proof}
控制前两条轨迹相同，故 $t_1,S^*$ 相同。退出时的二维状态均为 $(S_c,\eta)$，由常规系统解的唯一性，退出后的 $S,I$ 轨迹只相差时间平移。因此，两种策略从退出到 $I=1$ 所需的时间相同，退出后的新增感染积分也相同。分别消去两策略共同的前后两段，再用式\eqref{eq:joint:cumulative}，得到所列等式。其余仓室在退出时可以不同，故这里不推断隔离感染者或医疗占用也具有相同轨迹。
\end{proof}

在给定能力限制下，每一个处于
$[\Delta t[\overline c],\Delta t[\underline c]]$ 内的持续时间都可实现。事实上，取
$c_c(S)=(1-x)\underline c(S)+x\overline c(S)$、$x\in[0,1]$，式\eqref{eq:joint:duration}连续地连接两端；若两界并非几乎处处相等，则该持续时间随 $x$ 严格下降。结合推论\ref{cor:joint:whole-epidemic}可知，在固定阈值下，只对控制时长、清零时间和累计感染施加上限时，存在可行策略当且仅当最快控制 $\overline c$ 满足这些上限。再加入成本上限则不能仅用这一判据，因为最短时间与最小成本对应的控制未必相同。


规定退出状态还带来一个不依赖具体分配的感染下界。
\begin{proposition}
\label{prop:exit-infection-lower-bound}
对于模型\eqref{eq:model:full}，若 $0\le q(t)\le1$、$c(t)\ge0$，且某一有限时刻 $T$ 满足 $S(T)=S_c<S_0$，则
\begin{equation}
\Itcum(T)\ge\beta(S_0-S_c).
\label{eq:exit-infection-lower-bound}
\end{equation}
若整个 $[0,T]$ 上另有 $q(t)\le q_{\rm cap}<1$，则
\begin{equation}
\Itcum(T)\ge
\frac{\beta}{\beta+(1-\beta)q_{\rm cap}}(S_0-S_c).
\label{eq:exit-infection-capped}
\end{equation}
\end{proposition}
\begin{proof}
由易感者方程和累计新增感染的定义，
\[
\frac{d\Itcum}{dt}
=\frac{\beta}{\beta+(1-\beta)q(t)}(-\dot S).
\]
$-\dot S\ge0$，且比例系数分别不小于 $\beta$ 和
$\beta/[\beta+(1-\beta)q_{\rm cap}]$。在 $[0,T]$ 上积分即得两式。
\end{proof}
这一必要条件不保证下界总能达到，也不比较采用其他退出规则的方案。它表明，较低感染峰值、较少累计感染与规定退出状态可能不相容，不能仅通过重新分配两种干预消除这一限制。

\subsection{二次成本下的类内最优分配}
\label{sec:joint:cost}
持续时间最短的控制未必使式\eqref{eq:cost:J}最小，因为接触减少与追踪隔离的强度代价不同。由式\eqref{eq:joint:q}和状态时间变换，成本可直接写成接触率状态函数的积分：
\begin{equation}
J[c_c]=\frac{N}{\eta}\int_{S_c}^{S^*}f(S,c_c(S))\,dS,
\label{eq:joint:cost}
\end{equation}
其中
\begin{equation}
f(S,c)=
\frac{(1-c/c_0)^2+2\left(1-\dfrac{c_0S_c}{cS}\right)^2}
{cS-c_0\bar S}.
\label{eq:joint:cost-integrand}
\end{equation}
这里的 $c$ 是固定状态 $S$ 上待比较的接触率取值。分母来自易感者变化速度，因而要比较的是每减少一名社区易感者所对应的成本，不是仅使瞬时控制强度最小。

有能力限制时使用式\eqref{eq:joint:bounds}的上下界；无额外限制时取 $\underline c(S)=c_0S_c/S$、$\overline c(S)=c_0$。以下假定允许区间非空，且不增加总时长、调整次数或控制变化率等约束。
\begin{theorem}
\label{thm:joint:cost-minimum}
在规定的全部有界可测阈值控制中，二次成本最小值可达到，且
\begin{equation}
\min_{c_c}J[c_c]
=\frac{N}{\eta}\int_{S_c}^{S^*}
\min_{\underline c(S)\le c\le\overline c(S)}f(S,c)\,dS.
\label{eq:joint:minimum-cost}
\end{equation}
允许函数 $c_c^*(S)$ 达到该值，当且仅当它几乎处处属于相应的一维最小化解集。对于每个 $S>S_c$，只需比较 $\underline c(S),\overline c(S)$ 与区间内满足
\begin{equation}
\begin{split}
0={}&\frac{c^3}{c_0^2}(c-c_0)
\left(c+c_0-\frac{2c_0\bar S}{S}\right)\\
&+2\left(c-\frac{c_0S_c}{S}\right)
\left(-c^2+\frac{3c_0S_c}{S}c-
\frac{2c_0^2\bar S S_c}{S^2}\right)
\end{split}
\label{eq:joint:stationary}
\end{equation}
的全部实根处的 $f$ 值。该方程为五次多项式方程，故候选点至多七个；重复端点与重根只计一次。$S=S_c$ 时允许区间退化为 $c=c_0$，无需求根。
\end{theorem}
\begin{proof}
式\eqref{eq:joint:cost}的分母满足
$c_c(S)S-c_0\bar S\ge c_0(S_c-\bar S)>0$，所以 $f$ 在允许区域上连续有限。每个状态上的允许区间紧且非空，故最小值存在。任意允许控制的成本均不小于式\eqref{eq:joint:minimum-cost}右侧。

为证明该下界可达到，先在本证明中将允许区间写为
$c=\underline c(S)+x[\overline c(S)-\underline c(S)]$，$x\in[0,1]$，并记
\[
H(S,x)=f\!\left(S,\underline c(S)+x[\overline c(S)-\underline c(S)]\right),
\quad m(S)=\min_{0\le x\le1}H(S,x).
\]
$H$ 在紧集上连续，因而 $m$ 连续。取各最小化解集中最小的元素 $x_*(S)$。对于任意 $0<a\le1$，
\[
\{S:x_*(S)<a\}
=\bigcup_{\substack{r\in\mathbb Q\\0\le r<a}}
\left\{S:\min_{0\le x\le r}H(S,x)=m(S)\right\}.
\]
右侧为可数个闭集的并，故 $x_*$ 可测。于是
$c_c^*(S)=\underline c(S)+x_*(S)[\overline c(S)-\underline c(S)]$
是可测最小化选择，由定理\ref{thm:joint:representation}产生可行时间控制并达到下界。任意控制与逐点最小值之差非负，因此积分差为零当且仅当逐点差几乎处处为零。

固定 $S$，对式\eqref{eq:joint:cost-integrand}关于 $c$ 求偏导。将导数乘以正数
$c^3(cS-c_0\bar S)^2/S$，便得到式\eqref{eq:joint:stationary}右侧。该多项式最高次系数为 $1/c_0^2>0$，不恒为零，至多有五个实根。连续可微函数在闭区间上的最小值只能在端点或内部驻点取得，故比较所列候选值即可；不需要假设 $f$ 单峰或凸。
\end{proof}
$\beta=1$ 时 $\bar S=0$，约去非零因子 $c$ 后，驻点条件简化为
\begin{equation}
\frac{c^4}{c_0^2}-3c^2+\frac{8c_0S_c}{S}c
-6\left(\frac{c_0S_c}{S}\right)^2=0.
\label{eq:joint:quartic}
\end{equation}
一般五次判据不提供最优控制的初等时间表达，也不保证最小化选择唯一或连续。其作用是给出完整的有限候选判据。

显式族可作为成本比较的参照。仅减少接触时，
\begin{equation}
J=\frac{N}{c_0\eta(S_c-\bar S)}
\left[S^*-\frac{S_c^2}{S^*}-2S_c\ln\frac{S^*}{S_c}\right],
\label{eq:joint:contact-cost}
\end{equation}
仅加强隔离时恢复式\eqref{eq:cost:state}。全部允许控制中的最小成本不高于式\eqref{eq:joint:allocation-range}内任一显式策略的成本，但显式族未必包含类内最优解，不能以该子族中的最小值替代式\eqref{eq:joint:minimum-cost}。


同一状态上的最优分配不显含 $\eta$，这一点还能确定感染阈值对类内最小成本的影响。
\begin{proposition}
\label{prop:joint:minimum-threshold}
固定 $N,S_0,I_0,\beta,\gamma,c_0,q_0$ 及能力界限。在控制被触发且满足式\eqref{eq:joint:feasibility}的阈值范围内，记 $J_{\min}(\eta)$ 为式\eqref{eq:joint:minimum-cost}的最小值，则 $J_{\min}$ 随 $\eta$ 严格下降。对于任意两个可行阈值，可以选择类内最优控制，使其在共同经过的易感者状态上采用相同的 $c_c^*(S),q_c^*(S)$。

固定阈值时，降低 $c_{\min}$ 或提高 $q_{\rm cap}$ 不会增大类内最小成本。
\end{proposition}
\begin{proof}
在本证明中记
$m(S)=\min_{\underline c(S)\le c\le\overline c(S)}f(S,c)$。
被积函数及允许界限均不含 $\eta$，且 $m$ 连续。对于 $S>S_c$，两项强度不可能同时为零，故 $m(S)>0$。因此
\[
J_{\min}(\eta)=\frac{N}{\eta}\int_{S_c}^{S^*(\eta)}m(S)\,dS.
\]
由 $\partial S^*/\partial\eta<0$，在可行区间内部有
\begin{equation}
\frac{dJ_{\min}}{d\eta}
=-\frac{J_{\min}}{\eta}
+\frac{N}{\eta}m(S^*)\frac{\partial S^*}{\partial\eta}<0.
\label{eq:joint:mincost-eta}
\end{equation}
端点处的严格比较也可直接由两个积分的上下限和正系数得到。将较低阈值下的逐状态最小化选择限制到较高阈值所经过的区间，仍在每个状态达到同一 $m(S)$，故产生兼容的最优控制。最后，扩大任一实施能力使每个状态的允许区间包含原区间，其最小被积函数值不增，积分后即得最后的结论。
\end{proof}
提高阈值时，最优状态函数可以保持不变，但时间参数化与开始采用控制的状态仍会改变。上述成本结论也不意味着全过程感染人数随之减少。

二次成本还能够说明，为什么仅隔离控制虽达到最短时间，却一般不是成本最小的分配。
\begin{proposition}
\label{prop:joint:terminal-allocation}
采用式\eqref{eq:cost:J}的权重 $1,2$。无额外能力限制时，任一逐状态最小化选择在充分接近退出状态的 $S>S_c$ 上满足
\begin{equation}
\frac{c_0S_c}{S}<c_c^*(S)<c_0,
\qquad q_0<q_c^*(S)<1,
\label{eq:joint:terminal-interior}
\end{equation}
即两种干预均严格强于常规水平。而且
\begin{equation}
\begin{aligned}
\lim_{S\downarrow S_c}
\frac{(c_0-c_c^*(S))/c_0}{1-S_c/S}&=\frac23,\\
\lim_{S\downarrow S_c}
\frac{(q_c^*(S)-q_0)/(1-q_0)}{1-S_c/S}&=\frac13.
\end{aligned}
\label{eq:joint:terminal-ratio}
\end{equation}
因此，类内最小成本严格小于仅减少接触和仅加强隔离两种策略各自的成本。加入满足 $c_{\min}<c_0$、$q_{\rm cap}>q_0$ 的固定能力限制后，式\eqref{eq:joint:terminal-interior}--\eqref{eq:joint:terminal-ratio}仍成立；与某一种单控制的严格成本比较仅在该单控制整段可行时使用。
\end{proposition}
\begin{proof}
在本证明中令 $z=1-S_c/S$，并将允许接触率写为
$c=c_0(1-zx)$、$0\le x\le1$。于是
\[
\frac{q_c-q_0}{1-q_0}=\frac{z(1-x)}{1-zx},\qquad
\frac{f(S,c_0(1-zx))}{z^2}
=\frac{x^2+2(1-x)^2/(1-zx)^2}
{c_0[S(1-zx)-\bar S]}.
\]
当 $S\downarrow S_c$ 时，右侧在 $x\in[0,1]$ 上一致收敛于
\[
\frac{x^2+2(1-x)^2}{c_0(S_c-\bar S)},
\]
该函数的唯一最小化点为 $x=2/3$。由一致收敛和紧性，任何原问题最小化点序列均趋于 $2/3$：否则可取收敛子列，其极限将是极限函数的另一个最小化点，产生矛盾。因此所有最小化点在充分接近退出时均位于 $(0,1)$，并得到式\eqref{eq:joint:terminal-ratio}。

在这样一个正长度状态区间上，两种单控制分别对应 $x=0$ 和 $x=1$，均不能达到逐状态最小值；其非负成本差的积分因而严格为正。有能力限制且 $c_{\min}<c_0,q_{\rm cap}>q_0$ 时，充分接近 $S_c$ 的原允许区间仍为 $[c_0S_c/S,c_0]$，因此同一局部论证成立。
\end{proof}
该结论只刻画退出前的最优分配，不证明整个控制期的最小化选择唯一、连续或单调。若将二次权重改为 $w_c,w_q>0$，同一证明把式\eqref{eq:joint:terminal-ratio}中的 $2/3,1/3$ 分别替换为 $w_q/(w_c+w_q),w_c/(w_c+w_q)$。这比较的是两种归一化强度，不是接触率与隔离人数的直接比例。

上述逐状态分配利用了成本中没有控制导数或跨状态的积分限制。若另行限定控制调整速度或隔离人天等跨状态预算，不同状态处的选择会受到共同限制，不能直接沿用式\eqref{eq:joint:minimum-cost}。若改变隔离费用的计入方式，则需重新检查成本是否仍能写成逐状态分离的积分。若仅将权重 $1,2$ 改为给定正常数 $w_c,w_q$，则相应替换式\eqref{eq:joint:cost-integrand}中的权重，存在性证明不变；式\eqref{eq:joint:stationary}两项的系数分别改为 $w_c,w_q$。下面分别说明无额外能力限制时的端点局部性质，以及取得匹配乘子解时对控制时长限制的处理。

命题\ref{prop:joint:terminal-allocation}只刻画退出附近的分配。对于整个控制期，定义状态阈值
\begin{equation}
S_{\rm sw}=\frac{3S_c+\sqrt{9S_c^2-8S_c\bar S}}{2}.
\label{eq:joint:ssw}
\end{equation}
由 $0\le\bar S<S_c$，有 $2S_c<S_{\rm sw}\le3S_c$。这一阈值只用于以下端点的局部判据，不是一般权重下全局分配的切换位置。
\begin{proposition}
\label{prop:joint:global-structure}
设 $w_c,w_q>0$，且无额外能力限制。对每个 $S\in(S_c,S^*]$，仅减少接触的取值 $c=c_0S_c/S$ 不是 $f(S,\cdot)$ 在允许区间 $[c_0S_c/S,c_0]$ 上的局部最小点。仅隔离的取值 $c=c_0$ 在 $S>S_{\rm sw}$ 时是相对于该区间的严格局部最小点，在 $S<S_{\rm sw}$ 时不是局部最小点。

因此，任一类内最小成本控制在几乎所有 $S\in(S_c,S^*]$ 上满足 $q_c^*(S)>q_0$；在 $S_c<S<\min\{S_{\rm sw},S^*\}$ 上还满足 $c_0S_c/S<c_c^*(S)<c_0$，即两种措施同时强于常规水平。
\end{proposition}
\begin{proof}
在本证明中记 $n(c)=w_c(1-c/c_0)^2+w_q(1-c_0S_c/(cS))^2$。$f$ 关于 $c$ 的导数与
\[
n'(c)(cS-c_0\bar S)-S n(c)
\]
同号。在左端点 $c=c_0S_c/S$，该量为
\[
-2w_c\left(1-\frac{S_c}{S}\right)(S_c-\bar S)
-Sw_c\left(1-\frac{S_c}{S}\right)^2<0.
\]
故向允许区间内部增大 $c$ 即可降低成本。在右端点 $c=c_0$，该量为
\[
-\frac{w_q}{S}\left(1-\frac{S_c}{S}\right)
\bigl(S^2-3S_cS+2S_c\bar S\bigr).
\]
括号内二次式的较小根小于 $S_c$，较大根为 $S_{\rm sw}$，从而在 $S>S_{\rm sw}$ 时导数为负，在 $S_c<S<S_{\rm sw}$ 时导数为正。结合端点处的单侧比较，得到所述局部性质。由定理\ref{thm:joint:cost-minimum}，类内最优控制几乎处处取逐状态最小点；它不取左端点，在后一状态区间也不取右端点，故得到内部分配结论。
\end{proof}
当 $S=S_{\rm sw}$ 时上述一阶导数为零，此处不作局部分类。$S>S_{\rm sw}$ 时，右端点是否达到全局最小仍取决于权重及其他候选值，不能仅由这一局部判据确定。

控制持续时间限制可以在取得匹配的乘子解时处理，但这不是一般跨状态约束的无条件解法。
\begin{remarkn}
\label{rem:joint:duration-constraint}
固定同一允许控制类。对 $\kappa\in\mathbb R$，将 $f(S,c)$ 换为
\[
f(S,c)+\frac{\kappa}{cS-c_0\bar S},
\]
即对 $J+\kappa\Delta t$ 逐状态最小化，定理\ref{thm:joint:cost-minimum}的存在性论证仍成立。记所得允许控制为 $c_\kappa$，若 $\Delta t[c_\kappa]=T$，则 $\kappa\ge0$ 时它在约束 $\Delta t\le T$ 下使 $J$ 最小，$\kappa\le0$ 时它在约束 $\Delta t\ge T$ 下使 $J$ 最小；对于等式约束 $\Delta t=T$，任意符号的 $\kappa$ 均给出这一充分条件。事实上，对满足相应约束的任一允许控制 $c_c$，有
\[
J[c_c]\ge J[c_\kappa]+\kappa\bigl(\Delta t[c_\kappa]-\Delta t[c_c]\bigr)\ge J[c_\kappa].
\]
在 $0<\beta<1$ 时，由式\eqref{eq:joint:quarantine}，时长下限 $\Delta t\ge T$ 等价于控制期新增隔离易感者人数不超过 $(1-\beta)[S^*-S_c-\gamma\eta T]$。上述充分条件不要求可达集凸，但不保证每个允许时长都有匹配的 $\kappa$，也不直接解决控制变化率或其他累计预算约束。
\end{remarkn}


\FloatBarrier

